\begin{table}
\centering
\begin{small}
\begin{tabular}{|l|ccc|ccc|cc|c|}
\hline
platform & \multicolumn{3}{c|}{x86-64 CPU} & \multicolumn{3}{c|}{Arm64 CPU (M4)} & \multicolumn{2}{c|}{CUDA (T4)} & \multicolumn{1}{c|}{MPS (M4)} \\
library & NumPy & JAX & PyTorch & NumPy & JAX & PyTorch & JAX & PyTorch & PyTorch \\ \hline
atan2 & 0.501 & 0.501 & 0.501 & 0.501 & 0.501 & 0.501 & 0.501 & 0.501 & 0.501 \\
hypot & 0.501 & 1.95 & 0.501 & 0.501 & 1.92 & 0.501 & 1.95 & 0.501 & 0.501 \\
pow & 0.501 & 0.501 & 0.501 & 0.501 & 0.501 & 0.501 & 0.501 & 0.501 & 0.501 \\
\hline
\end{tabular}
\end{small}
\caption{binary16 bivariate: largest error over all pairs (exhaustive). Values: largest error in ulps under the rules of Gladman et al. (3 significant digits rounded up; an input where only one of the library result and the exact result is NaN gives Inf). In parentheses, when different: the largest error excluding results affected by a flush-to-zero/denormals-are-zero mode of the backend and NaN mismatches. 0.500 = correctly rounded on all inputs; -- = function not provided for this format/backend; marks: f = some results flushed by the backend's FTZ/DAZ mode, n = spurious NaN for some inputs, m = NaN missing (number returned where the result is NaN). }
\label{tab:binary16bi}
\end{table}
